% =============================================================================
% MLSys·im Tutorial Tutorial — Parts 0–4 (Morning Session)
% =============================================================================
\documentclass[aspectratio=169, 12pt]{beamer}
\usepackage{../../../slides/assets/beamerthememlsys}

\mlsyssetup{
  volume       = {Tutorial},
  chapter      = {Module 3},
  logo         = {../../../slides/assets/img/logo-mlsysbook.png},
  instlogo     = {../../../slides/assets/img/logo-harvard.png},
  chaptertitle = {MLSys·im: Scale, Dollars, and Carbon},
}

% --- Fonts ---
\usepackage[T1]{fontenc}
\usepackage[scaled=0.9]{helvet}
\usepackage{courier}
\renewcommand{\familydefault}{\sfdefault}

% --- Packages ---
\usepackage{amsmath}
\usepackage{booktabs}
\usepackage[table]{xcolor}
\usepackage{listings}
\usepackage{tikz}
\usetikzlibrary{arrows.meta, positioning, calc, decorations.pathreplacing}

% --- Code listings ---
\lstset{
  language=Python,
  basicstyle=\ttfamily\footnotesize,
  keywordstyle=\color{crimson}\bfseries,
  stringstyle=\color{datastroke},
  commentstyle=\color{midgray}\itshape,
  backgroundcolor=\color{computeblue!20},
  frame=single,
  rulecolor=\color{computestroke},
  numbers=none,
  breaklines=true,
  columns=fullflexible,
  keepspaces=true,
  showstringspaces=false,
  xleftmargin=4pt,
  xrightmargin=4pt,
  aboveskip=6pt,
  belowskip=4pt,
}

% --- Convenience macros ---
\newcommand{\mlsysim}{\texttt{mlsysim}}
\newcommand{\wallbox}[2]{%
  \begin{block}{#1}#2\end{block}%
}
\newcommand{\PredictStart}{\begin{alertblock}{Predict Before You Peek}}
\newcommand{\PredictEnd}{\end{alertblock}}

% --- Image paths ---
\graphicspath{{images/}}

% --- Section count (must match actual \section{} count) ---
\setcounter{mlsystotalsections}{6}

\title{MLSys·im Tutorial --- Module 3}
\subtitle{Scale, Dollars, and Carbon}
\author{Vijay Janapa Reddi}
\institute{Harvard University}
\date{Tutorial}

% =============================================================================
\begin{document}

% --- Roadmap ---
\begin{frame}{Roadmap: Conference Tutorial}
\centering\small
\begin{tabular}{rll}
\toprule
\textbf{Module} & \textbf{Topic} & \textbf{Status} \\
\midrule
Module 1 & Foundations \& Architecture & \checkmark Done \\
Module 2 & Advanced Single-Node Analysis & \checkmark Done \\
\rowcolor{crimson!12}
\textbf{Module 3} & \textbf{Scale, Dollars, and Carbon} & \textbf{$\leftarrow$ You are here} \\
Module 4 & Design Space Exploration \& Synthesis & \\
\bottomrule
\end{tabular}
\end{frame}


\section{Going Distributed}

% --- Slide 4.1: Key Question ---
\begin{frame}{Key Question}
\note{[1 min] ``Your model does not fit on one GPU. Or it fits but
training would take a year. Either way, you need more GPUs.
But adding GPUs is not free.''
% -- FLEX: [CORE]
}

\centering
\Large\bfseries
If 1 GPU takes 30 days,\\[0.2cm]
do 1000 GPUs take 43 minutes?
\end{frame}

\begin{frame}{The Algebra of Distributed Overheads}
\begin{block}{Communication Time (e.g., Ring AllReduce)}
\[
T_{\text{comm}} = \alpha + \frac{\text{Message Size}}{\text{Network Bandwidth}}
\]
\end{block}
\textbf{The Physics:}
\begin{itemize}
  \item $\alpha$ (Network Latency) dominates for small messages.
  \item Bandwidth dominates for large gradient synchronizations.
\end{itemize}
\vfill
\begin{alertblock}{The Tool}
The \texttt{DistributedModel} solver automatically runs this algebra across the specified network fabric (e.g., NVLink vs. InfiniBand) to compute exactly how much time GPUs sit idle.
\end{alertblock}
\end{frame}


% --- Slide 4.2: Why Distribute? ---
\begin{frame}{Why Distribute?}
\note{[2 min] ``Two reasons to go distributed: (1) the model does not
fit in one GPU's memory, or (2) you want to finish sooner.
Reason 1 is a hard constraint. Reason 2 is an optimization.''
% -- FLEX: [CORE]
}

\small
\textbf{Reason 1: Model does not fit}
\begin{itemize}
  \item Llama-3 70B FP16 $=$ 140 GB $>$ H100's 80 GB
  \item \alert{Must} split across at least 2 GPUs
\end{itemize}

\vspace{0.3cm}
\textbf{Reason 2: Time-to-train}
\begin{itemize}
  \item 1 H100 training Llama-3 70B $\approx$ 15 GPU-years
  \item 1024 H100s $\approx$ 5 days (if scaling were perfect)
  \item But scaling is \textit{never} perfect...
\end{itemize}

\vspace{0.3cm}
\centering
\begin{tikzpicture}[scale=0.7, >=Stealth]
  \draw[->, thick] (0,0) -- (7,0) node[right, font=\scriptsize] {GPUs};
  \draw[->, thick] (0,0) -- (0,4) node[above, font=\scriptsize] {Speedup};
  \draw[dashed, midgray] (0,0) -- (6.5,3.9) node[right, font=\scriptsize\itshape] {ideal};
  \draw[very thick, crimson] (0,0) .. controls (2,2) and (4,3) .. (6.5,3.2);
  \node[font=\scriptsize, crimson] at (5.5, 2.4) {reality};
\end{tikzpicture}
\end{frame}

% --- Slide 4.3: The Three Dimensions of Parallelism ---
\begin{frame}{3D Parallelism: DP $\times$ TP $\times$ PP}
\note{[3 min] ``Every distributed strategy is a combination of three
dimensions.'' WARN: Students often confuse TP and PP.
TP splits within a layer; PP splits between layers.
% -- FLEX: [CORE]
}

\small
\begin{columns}[T]
\begin{column}{0.32\textwidth}
\textbf{Data Parallel (DP)}
\begin{itemize}\setlength\itemsep{1pt}\footnotesize
  \item Replicate full model
  \item Split data across replicas
  \item AllReduce gradients
  \item \textit{Most common}
\end{itemize}
\end{column}
\begin{column}{0.32\textwidth}
\textbf{Tensor Parallel (TP)}
\begin{itemize}\setlength\itemsep{1pt}\footnotesize
  \item Split each layer's weights
  \item Split activations, not data
  \item AllReduce per layer (2$\times$!)
  \item Needs fast interconnect
\end{itemize}
\end{column}
\begin{column}{0.32\textwidth}
\textbf{Pipeline Parallel (PP)}
\begin{itemize}\setlength\itemsep{1pt}\footnotesize
  \item Split model into stages
  \item Each GPU owns $L/\text{PP}$ layers
  \item Pipeline bubbles
  \item Needs less bandwidth
\end{itemize}
\end{column}
\end{columns}

\vspace{0.3cm}
Total GPUs: $N = \text{DP} \times \text{TP} \times \text{PP}$

\vspace{0.2cm}
\centering
\scriptsize
\begin{tabular}{lccc}
\toprule
\textbf{Property} & \textbf{DP} & \textbf{TP} & \textbf{PP} \\
\midrule
Splits & Data & Weights + Activations & Layers \\
Communication & AllReduce (gradients) & AllReduce (activations) & Point-to-point \\
BW requirement & Moderate & Very high (NVLink) & Low \\
Bubble overhead & None & None & $\sim(P{-}1)/(M{+}P{-}1)$ \\
\bottomrule
\end{tabular}
\end{frame}

% --- Slide 4.3b: AllReduce Concrete Example ---
\begin{frame}[fragile]{AllReduce: A Concrete Example}
\note{[2 min] NUMBERS FIRST, then formula. Students need to see the magnitude before the algebra.}

\small
\textbf{Setup:} 8 H100 GPUs, NVLink at 900 GB/s, Llama-3 8B (16 GB gradients)

\vspace{0.3cm}
\begin{enumerate}
\item Each GPU computes its local gradient: \textbf{16 GB}
\item All 8 GPUs must end up with the \textbf{same averaged gradient}
\item Ring AllReduce passes chunks around the ring\ldots
\end{enumerate}

\vspace{0.3cm}
\begin{lstlisting}
t = mlsysim.physics.calc_ring_allreduce_time(
    message_bytes=16e9,
    n_gpus=8,
    bandwidth_bytes_s=900e9,
    latency_s=500e-9,
)
print(f"AllReduce time: {t.to('ms'):.1f}")
# -> ~35 ms (bandwidth-dominated, latency is negligible)
\end{lstlisting}

\vfill
\centering
\textbf{35 ms} to synchronize 8 GPUs. Now: what happens at 256 GPUs?
\end{frame}

% --- Slide 4.4: Data Parallelism + AllReduce ---
\begin{frame}{Wall 14: The Communication Wall (AllReduce)}
\note{[3 min] ``Quick: 1 GB of gradients, 8 GPUs, 50 GB/s NVLink.
How long for AllReduce?'' Expected: 2*(7/8)*1/50 = 35 ms.
Ring AllReduce sends 2(N-1)/N times the data. As N grows,
this approaches 2x.
% -- FLEX: [CORE]
}

\small
\wallbox{The Communication Wall}{
\[
  T_{\text{AllReduce}} = \frac{2(N-1)}{N} \times \frac{M}{BW}
    + 2(N-1) \times \text{latency}
\]
}

\vspace{0.2cm}
\textbf{Example:} 1 GB gradients, 8$\times$ H100 on NVLink (900 GB/s)

\[
T = \frac{2 \times 7}{8} \times \frac{1}{900} \approx 1.9\;\text{ms}
\]

\vspace{0.2cm}
\textbf{Same gradients}, 256$\times$ H100 across InfiniBand (50 GB/s):

\[
T = \frac{2 \times 255}{256} \times \frac{1}{50} \approx 40\;\text{ms}
\]

\vspace{0.2cm}
\alert{NVLink is 18$\times$ faster than InfiniBand for AllReduce.\\
That is why TP must stay within a node.}
\end{frame}

% --- Slide 4.5: Tensor Parallelism ---
\begin{frame}{Tensor Parallelism: Splitting Layers}
\note{[3 min] ``TP splits each layer's weight matrix across GPUs.
Every forward and backward pass requires 2 AllReduce ops per
layer. That is why TP only works on NVLink, not across nodes.''
% -- FLEX: [CORE]
}

\small
\textbf{How it works:}
\begin{enumerate}\setlength\itemsep{2pt}
  \item Split weight matrix $W$ column-wise across $T$ GPUs
  \item Each GPU computes $Y_i = X \cdot W_i$ (partial result)
  \item AllReduce to combine: $Y = \sum Y_i$
  \item \alert{2 AllReduce ops per layer} (forward + backward)
\end{enumerate}

\vspace{0.3cm}
\textbf{TP overhead:}
\[
  T_{\text{TP}} = 2 \times L \times T_{\text{AllReduce}}(T)
\]

\begin{exampleblock}{Llama-3 70B, TP=8 on NVLink (900 GB/s)}
  \begin{itemize}
    \item 80 layers $\times$ 2 AllReduce $\times$ $\sim$0.1 ms each $\approx$ \textbf{16 ms overhead per step}
    \item This is 10--20\% of a typical training step
  \end{itemize}
\end{exampleblock}
\end{frame}

% --- Slide 4.6: Pipeline Parallelism ---
\begin{frame}{Pipeline Parallelism: The Bubble Problem}
\note{[3 min] ``With 4 stages and 4 microbatches, what fraction
of time is wasted?'' Expected: 3/7 = 43\%. With 32 microbatches:
3/35 = 8.6\%. Lesson: more microbatches = smaller bubble.
% -- FLEX: [CORE]
}

\small
\wallbox{Pipeline Bubble Fraction}{
\[
  \text{Bubble} = \frac{P - 1}{M + P - 1}
\]
where $P$ = pipeline stages, $M$ = microbatches
}

\vspace{0.3cm}
\scriptsize
\begin{tabular}{lcccc}
\toprule
$P$ (stages) & $M$ (microbatches) & Bubble & Effective utilization \\
\midrule
4 & 4 & 43\% & 57\% \\
4 & 16 & 16\% & 84\% \\
4 & 32 & 8.6\% & 91\% \\
8 & 32 & 18\% & 82\% \\
\bottomrule
\end{tabular}

\normalsize
\vspace{0.2cm}
\alert{More microbatches $\Rightarrow$ smaller bubble.\\
But more microbatches = more memory for activations.}
\end{frame}

% --- Slide 4.7: Gradient Accumulation ---
\begin{frame}{Gradient Accumulation: Virtual Batch Size}
\note{[2 min] ``Process K small microbatches and accumulate gradients
before the optimizer step. This fills the pipeline and amortizes AllReduce.''
% -- FLEX: [OPTIONAL]
}

\small
\[
  B_{\text{effective}} = B_{\text{micro}} \times K \times \text{DP}
\]

\textbf{Why accumulate?}
\begin{itemize}\setlength\itemsep{2pt}
  \item Fill the pipeline ($M = K$ microbatches)
  \item Amortize AllReduce cost over $K$ steps
  \item Simulate large batch size without large memory
  \item Trade compute (more forward passes) for communication (fewer AllReduce)
\end{itemize}

\vspace{0.2cm}
\textbf{Example:} DP=128, $B_{\text{micro}}$=4, $K$=8
\begin{itemize}
  \item $B_{\text{effective}} = 4 \times 8 \times 128 = 4096$
  \item AllReduce only once per 8 microbatches
  \item Pipeline bubble: $(P-1)/(8+P-1)$ --- much smaller
\end{itemize}
\end{frame}

% --- Slide 4.8: Hierarchical Communication ---
\begin{frame}{Hierarchical AllReduce: NVLink + InfiniBand}
\note{[2 min] ``Hierarchical AllReduce first reduces within each
node (fast NVLink), then across nodes (slower IB), then
broadcasts back. This exploits the bandwidth hierarchy.''
% -- FLEX: [OPTIONAL]
}

\small
\textbf{Real cluster topology:}

\centering
\begin{tikzpicture}[scale=0.75, >=Stealth,
  gpu/.style={draw, fill=computeblue, rounded corners, minimum width=0.6cm,
              minimum height=0.5cm, font=\tiny},
  node/.style={draw, fill=white, rounded corners=4pt, dashed, inner sep=4pt}]

  % Node 0
  \node[node, label=above:{\scriptsize Node 0}] (n0) at (0,0) {
    \begin{tikzpicture}
      \foreach \i in {0,...,3} {
        \node[gpu] (g0\i) at (\i*0.8, 0) {G\i};
      }
    \end{tikzpicture}
  };

  % Node 1
  \node[node, label=above:{\scriptsize Node 1}] (n1) at (5.5,0) {
    \begin{tikzpicture}
      \foreach \i in {0,...,3} {
        \node[gpu] (g1\i) at (\i*0.8, 0) {G\i};
      }
    \end{tikzpicture}
  };

  % NVLink labels
  \node[font=\tiny, datastroke] at (0, -0.9) {NVLink 900 GB/s};
  \node[font=\tiny, datastroke] at (5.5, -0.9) {NVLink 900 GB/s};

  % IB link
  \draw[very thick, crimson, <->] (2.2, 0) -- (3.3, 0)
    node[midway, above, font=\tiny] {IB 50 GB/s};
\end{tikzpicture}

\vspace{0.3cm}
\flushleft
\small
\textbf{3-step hierarchical AllReduce:}
\begin{enumerate}\setlength\itemsep{1pt}
  \item \textbf{Local reduce} within each node (NVLink --- fast)
  \item \textbf{Global AllReduce} across leader GPUs (InfiniBand --- slow)
  \item \textbf{Local broadcast} within each node (NVLink --- fast)
\end{enumerate}

\alert{TP within node (NVLink). DP across nodes (InfiniBand).}
\end{frame}

% --- Slide 4.9: Live Demo --- DistributedModel ---
\begin{frame}[fragile]{Live Demo: Distributed Training Analysis}
\note{[3 min] Run DistributedModel live. Show communication overhead,
bubble fraction, and scaling efficiency.
% -- FLEX: [CORE]
}

\small
\begin{lstlisting}
fleet = mlsysim.Systems.Clusters.Frontier_8K
dist  = mlsysim.DistributedModel()
result = dist.solve(llama, fleet, batch_size=4096,
         tp_size=8, pp_size=1, microbatch_count=32,
         seq_len=4096)
print(f"Scaling eff:   {result.scaling_efficiency:.1%}")
print(f"Comm overhead: {result.communication_overhead:.1%}")
print(f"Effective MFU: {result.effective_mfu:.1%}")
\end{lstlisting}

\vspace{0.2cm}
\begin{exampleblock}{What to look for}
  Communication overhead + bubble fraction = total efficiency loss.
  Effective MFU $=$ single-node MFU $\times$ scaling efficiency.
\end{exampleblock}
\end{frame}

% --- Slide 4.10: Wall 15 --- The Fragility Wall ---
\begin{frame}{Wall 15: The Fragility Wall (Reliability)}
\note{[2 min] ``If you have 10,000 GPUs each with 50,000 hour MTBF,
what is the cluster MTBF?'' Expected: 50,000/10,000 = 5 hours.
This is why checkpointing exists.
% -- FLEX: [CORE]
}

\small
\wallbox{The Fragility Wall}{
\[
  \text{Cluster MTBF} = \frac{\text{Component MTBF}}{N_{\text{components}}}
\]
}

\vspace{0.2cm}
\begin{tabular}{lcc}
\toprule
\textbf{Scale} & \textbf{GPUs} & \textbf{Cluster MTBF} \\
\midrule
Research lab   & 8     & 260 days \\
Mid cluster    & 256   & 8 days \\
Large cluster  & 1,024 & 2 days \\
Frontier-scale & 8,192 & 6 hours \\
Mega cluster   & 100K  & 30 minutes \\
\bottomrule
\end{tabular}

\vspace{0.2cm}
\alert{At frontier scale, something breaks every 6 hours.\\
Without checkpointing, every failure wastes the entire run since the last save.}
\end{frame}

% --- Slide 4.10b: Predict --- Scaling to 256 GPUs ---
\begin{frame}{Predict: Scaling to 256 GPUs}
\note{[2 min] PREDICTION. Hands up for each answer. Most will say 256x.
Turn to your neighbor: did you get the same answer? Why or why not? 60 seconds.}
\centering
\Large
\textbf{You have 8 H100s doing data-parallel training.}\\[0.5cm]
\textbf{You scale to 256 GPUs.}\\[0.5cm]
\normalsize
How much faster will training be?\\[0.3cm]
\begin{enumerate}[(A)]
  \item 32$\times$ faster (perfect scaling)
  \item 20--25$\times$ faster
  \item 10--15$\times$ faster
  \item \textbf{It depends on the model size}
\end{enumerate}
\end{frame}

% --- Slide 4.11: Scaling Efficiency ---
\begin{frame}{Scaling Efficiency: The Amdahl Trap}
\note{[2 min] ``Scaling efficiency is the fraction of ideal speedup
you actually achieve.'' Includes comm overhead, pipeline
bubbles, stragglers, and failure recovery.
% -- FLEX: [CORE]
}

\small
\[
  \eta_{\text{scaling}} = \frac{\text{Actual speedup}}{N}
  = \frac{1}{1 + \text{comm\_frac} + \text{bubble\_frac} + \text{straggler\_frac}}
\]

\vspace{0.3cm}
\begin{columns}[T]
\begin{column}{0.55\textwidth}
\textbf{What eats scaling efficiency:}
\begin{enumerate}\setlength\itemsep{2pt}
  \item AllReduce communication
  \item Pipeline bubbles
  \item Straggler effects (slowest GPU)
  \item Checkpoint I/O
  \item Failure recovery
\end{enumerate}
\end{column}
\begin{column}{0.42\textwidth}
\centering
\scriptsize
\begin{tabular}{lc}
\toprule
\textbf{System} & $\eta_{\text{scaling}}$ \\
\midrule
8 GPUs (NVLink)  & 95--98\% \\
64 GPUs (IB)     & 85--92\% \\
1024 GPUs        & 70--85\% \\
8192 GPUs        & 55--70\% \\
\bottomrule
\end{tabular}
\end{column}
\end{columns}

\vspace{0.3cm}
\alert{At 8192 GPUs, you lose 30--45\% of your compute to overhead.}
\end{frame}

% --- Slide 4.12: Predict Before You Peek #5 ---
\begin{frame}{Predict: Optimal Parallelism Config}
\note{[2 min] You have 64 H100s. Llama-3 70B (140 GB FP16).
What TP x PP x DP? Give 90 seconds.
Expected: TP=8 (NVLink), PP=1 (no bubbles), DP=8 (64/8).
Turn to your neighbor: did you get the same answer? Why or why not? 60 seconds.
% -- FLEX: [CORE]
}

\centering
\Large\bfseries
64 H100s. Llama-3 70B (140 GB FP16).\\[0.3cm]
\normalsize
What is the optimal TP $\times$ PP $\times$ DP?\\[0.5cm]

\pause

\small
\begin{tabular}{lcccl}
\toprule
\textbf{Config} & \textbf{TP} & \textbf{PP} & \textbf{DP} & \textbf{Why} \\
\midrule
Candidate A & 8 & 1 & 8 & TP within node, no bubbles \\
Candidate B & 4 & 2 & 8 & Less TP comm, but has bubbles \\
Candidate C & 2 & 4 & 8 & Minimal TP, large bubble \\
\bottomrule
\end{tabular}

\vspace{0.3cm}
\pause
\alert{Candidate A is typically best:} TP=8 uses full NVLink bandwidth,
PP=1 avoids pipeline bubbles entirely, DP=8 across nodes.

\vspace{0.2cm}
\textit{Rule of thumb: maximize TP within a node, minimize PP.}
\end{frame}

% --- Slide 4.13: Exercise 4 ---
\begin{frame}[fragile]{Exercise 4: Distributed Training Design}
\note{[5 min] Sweep TP in [1,2,4,8] and PP in [1,2,4,8] for
Llama-3 70B on 64 H100s. Expected: TP=8, PP=1, DP=8.
Turn to your neighbor: did you get the same answer? Why or why not? 60 seconds.
% -- FLEX: [CORE]
}

\centering
\Large\bfseries
Hands-On Exercise\\[0.5cm]
\normalsize

\textbf{Question:} Find the optimal TP$\times$PP for Llama-3 70B on 64 H100s.

\vspace{0.3cm}
\begin{lstlisting}
llama70 = mlsysim.Models.Language.Llama3_70B
fleet   = mlsysim.Systems.Clusters.Research_256
for tp in [1, 2, 4, 8]:
    for pp in [1, 2, 4, 8]:
        if tp * pp > 64: continue
        r = dist.solve(llama70, fleet, batch_size=512,
            tp_size=tp, pp_size=pp, seq_len=4096,
            microbatch_count=max(4, 64//(tp*pp)))
        print(f"TP={tp} PP={pp} eff={r.scaling_efficiency:.1%}")
\end{lstlisting}
\end{frame}

% --- Slide 4.14: Straggler Effects ---
\begin{frame}{Stragglers: The Slowest GPU Sets the Pace}
\note{[2 min] ``In synchronous training, every GPU must finish
before the next step begins. At 1000 GPUs, even 1\% variation
means 10 GPUs are significantly slow on any given step.''
% -- FLEX: [OPTIONAL]
}

\small
\textbf{Synchronous training:} step time $=$ $\max_i(T_i)$

\vspace{0.2cm}
\begin{itemize}\setlength\itemsep{2pt}
  \item \textbf{Thermal throttling:} hot GPUs clock down 5--10\%
  \item \textbf{Network congestion:} some AllReduce messages delayed
  \item \textbf{OS jitter:} background tasks steal cycles
  \item \textbf{Memory pressure:} GC pauses in the data pipeline
\end{itemize}

\vspace{0.3cm}
\textbf{Mitigation strategies:}
\begin{itemize}\setlength\itemsep{2pt}
  \item Asynchronous SGD (trade accuracy for speed)
  \item Backup workers (redundant computation)
  \item Bounded staleness (allow slight divergence)
  \item \texttt{DistributedModel(straggler\_factor=1.05)} to simulate 5\% drag
\end{itemize}
\end{frame}

% --- Slide 4.15: Part 4 Key Takeaway ---
\begin{frame}{Part 4: Key Takeaway}
\note{[1 min] One sentence. Repeat.
``Distributed training is a communication problem disguised
as a compute problem.''
% -- FLEX: [CORE]
}

\centering\Large

\textbf{Distributed training is a communication problem\\
disguised as a compute problem.}\\[0.5cm]

\normalsize
\begin{itemize}
  \item 3D parallelism (DP $\times$ TP $\times$ PP) decomposes the problem.
  \item TP needs NVLink (within node). DP works over InfiniBand (across nodes).
  \item Pipeline bubbles shrink with more microbatches.
  \item Reliability degrades as $\text{MTBF}/N$ --- checkpointing is mandatory.
  \item \texttt{DistributedModel.solve()} captures all these effects.
\end{itemize}

\vspace{0.5cm}
\centering
\textit{Lunch break --- reconvene at 1:00 PM for Part 5.}
\end{frame}


\section{Economics \& Sustainability}

% --- 5.1 Key Question ---
\begin{frame}{Key Question}
\note{[1 min] Open with the dollar question. Pause dramatically.}
\begin{center}
\Large\bfseries
How much does it \emph{really} cost to train and serve an LLM---\\[4pt]
in dollars, kilowatt-hours, and tonnes of CO$_2$?
\end{center}
\vfill
\begin{quotation}
\small\itshape
``A 1024-GPU H100 cluster costs \$30M in hardware. Electricity is
surprisingly small ($\sim$10\% of total). The dominant cost lever is
utilization---idle GPUs cost the same as busy ones.''
\end{quotation}
\end{frame}

% --- 5.2 The TCO Equation ---
\begin{frame}[fragile]{Wall 17: The Capital Wall (TCO)}
\note{[3 min] Run the TCO code live. Key surprise: GPUs are only 40\% of total cost.}
\begin{columns}[T]
\column{0.55\textwidth}
\begin{block}{The Equation}
\[
\text{TCO} = \underbrace{\text{CapEx}}_{\text{hw + facility}}
          + \underbrace{\text{OpEx}_{\text{energy}}}_{\text{electricity}}
          + \underbrace{\text{OpEx}_{\text{maint}}}_{\text{staff}}
\]
\end{block}
\vspace{0.3em}
\small
\textbf{Key insight:} GPUs are only $\sim$40\% of total CapEx.\\
Networking, cooling, facility, and staff add 50--150\%.

\column{0.42\textwidth}
\begin{lstlisting}
econ = mlsysim.EconomicsModel()
r = econ.solve(
  fleet=mlsysim.Systems
    .Clusters.Research_256,
  duration_days=30,
  infrastructure_multiplier=2.5)
print(f"TCO: ${r.tco_usd:,.0f}")
\end{lstlisting}
\end{columns}

\vspace{0.3em}
\small
\begin{itemize}\setlength\itemsep{1pt}
  \item \texttt{infrastructure\_multiplier = 2.0--2.5} for full datacenter TCO
  \item Amortized over 3--5 year depreciation schedule
  \item Cloud spot instances can beat on-prem for bursty workloads
\end{itemize}
\end{frame}

% --- 5.3 Infrastructure Multiplier ---
\begin{frame}{The Infrastructure Multiplier}
\note{[2 min] Walk through the cost breakdown table. Emphasize: budgeting only for GPUs underestimates by 2-3x.}
\centering
\begin{tabular}{lrr}
\toprule
\textbf{Component} & \textbf{Cost (\%)} & \textbf{Multiplier} \\
\midrule
GPU Hardware       & 40\% & 1.0$\times$ \\
Network (IB)       & 20\% & --- \\
Power \& Cooling   & 15\% & --- \\
Facility / Land    & 10\% & --- \\
Staff / Operations & 15\% & --- \\
\midrule
\textbf{Total}     & 100\% & \textbf{2.5$\times$} \\
\bottomrule
\end{tabular}

\vspace{1em}
\alert{Pitfall:} Budgeting only for GPUs underestimates TCO by 2--3$\times$.
\end{frame}

% --- 5.4 The Sustainability Equation ---
\begin{frame}[fragile]{Wall 18: The Sustainability Wall}
\note{[3 min] Three levers: energy, PUE, carbon intensity. Geography dominates.}
\begin{columns}[T]
\column{0.55\textwidth}
\begin{block}{The Equation}
\small
\[
\text{CO}_2 = \text{Energy} \times \text{PUE} \times \text{CI}
\]
\end{block}
\vspace{0.3em}
\small
\textbf{Three levers:}
\begin{enumerate}\setlength\itemsep{1pt}
  \item \textbf{Energy} --- Better MFU, smaller model
  \item \textbf{PUE} --- Liquid cooling (1.3 $\to$ 1.06)
  \item \textbf{Carbon intensity} --- Choose your grid
\end{enumerate}

\column{0.42\textwidth}
\begin{lstlisting}
sus = mlsysim.SustainabilityModel()
r = sus.solve(
  fleet=mlsysim.Systems
    .Clusters.Research_256,
  datacenter=mlsysim.Infrastructure
    .Grids.Quebec,
  duration_days=30, mfu=0.45)
print(f"{r.carbon_footprint_kg
         /1000:.1f} tonnes CO2")
\end{lstlisting}
\end{columns}

\vspace{0.3em}
\begin{center}
\alert{Geography is the single biggest lever for sustainable AI.}
\end{center}
\end{frame}

% --- 5.5 Geography Matters ---
\begin{frame}{Carbon Intensity: The 41$\times$ Gap}
\note{[2 min] Show the table. Poland vs Quebec: 41x difference. Same energy, different grid.}
\begin{columns}[T]
\column{0.45\textwidth}
\centering
\begin{tabular}{lcr}
\toprule
\textbf{Region} & \textbf{Mix} & \textbf{gCO$_2$} \\
\midrule
Quebec    & Hydro       & 20 \\
Sweden    & Hydro+Nuc   & 25 \\
US Avg    & Mixed       & 390 \\
Germany   & Coal+Wind   & 380 \\
Poland    & Coal        & 820 \\
\bottomrule
\end{tabular}

\vspace{0.5em}
\small
Training the \emph{same model}:\\
\textbf{Poland:} $\sim$800\,t\quad vs\quad \textbf{Quebec:} $\sim$20\,t\\[4pt]
$\Rightarrow$ \textbf{41$\times$} difference.

\column{0.52\textwidth}
\centering
\includegraphics[width=0.85\textwidth]{images/pdf/carbon-geography.pdf}
\end{columns}
\end{frame}

% --- 5.6 Embodied Carbon ---
\begin{frame}{Embodied Carbon: Manufacturing Dominates at the Edge}
\note{[2 min] Key insight: at cloud scale, operational carbon dominates. At TinyML scale, embodied carbon dominates.}
\begin{columns}[T]
\column{0.5\textwidth}
\textbf{Cloud / Training}
\begin{itemize}
  \item Operational carbon dominates
  \item GPU runs 24/7 for months
  \item Embodied: $\sim$5\% of lifecycle
\end{itemize}

\column{0.5\textwidth}
\textbf{IoT / TinyML}
\begin{itemize}
  \item Manufacturing carbon dominates
  \item MCU uses microwatts
  \item Embodied: $>$90\% of lifecycle
  \item Multiplied by millions of devices
\end{itemize}
\end{columns}

\vspace{1em}
\begin{center}
\small
Source: Gupta et al.\ (2022), ``ACT: Designing Sustainable Computer\\
Systems with an Architectural Carbon Modeling Tool''
\end{center}
\end{frame}

% --- 5.7 Live Demo: Economics ---
\begin{frame}[fragile]{Live Demo: TCO + Carbon Analysis}
\note{[3 min] Run live. Show TCO and carbon side by side.}
\begin{lstlisting}
fleet = mlsysim.Systems.Clusters.Research_256
econ  = mlsysim.EconomicsModel()
result = econ.solve(
    fleet=fleet, duration_days=30,
    datacenter=mlsysim.Infrastructure.Grids.Quebec,
    mfu=0.45, infrastructure_multiplier=2.5,
    amortization_years=3.0)
print(f"TCO: ${result.tco_usd:,.0f}")
print(f"Carbon: {result.carbon_footprint_kg/1000:.1f} t")
\end{lstlisting}
\end{frame}

% --- 5.8 Live Demo: Geography Comparison ---
\begin{frame}[fragile]{Live Demo: The Carbon Geography Experiment}
\note{[2 min] Run live. Same energy, 41x less carbon. Geography wins.}
\begin{lstlisting}
fleet  = mlsysim.Systems.Clusters.Research_256
solver = mlsysim.SustainabilityModel()
for region in [mlsysim.Infrastructure.Grids.Poland,
               mlsysim.Infrastructure.Grids.Quebec]:
    r = solver.solve(fleet=fleet, duration_days=30,
                     datacenter=region, mfu=0.45)
    print(f"{r.region_name:>12}: "
          f"{r.carbon_footprint_kg/1000:.1f} t CO2")
\end{lstlisting}

\vspace{0.5em}
\begin{exampleblock}{Expected Output}
\ttfamily\scriptsize
\hspace{1em}Poland: 1,064 MWh, 872.5 tonnes CO2\\
\hspace{1em}Quebec: 1,064 MWh, \phantom{0}21.3 tonnes CO2
\end{exampleblock}

\vspace{0.3em}
\alert{Same energy. 41$\times$ less carbon. Geography wins.}
\end{frame}

% --- 5.9 Exercise Part A: Carbon-Aware Placement ---
\begin{frame}[fragile]{Exercise 5a: Carbon-Aware Placement}
\note{[3 min] Expected answer: $\sim$850 tonnes CO$_2$ saved.
The energy is identical ($\sim$2{,}100 MWh)---only carbon intensity differs.
Key insight: the cheapest carbon reduction is a \texttt{git push} to a different region.
Turn to your neighbor: did you get the same answer? Why or why not? 60 seconds.}

\begin{alertblock}{Your Task (3 minutes)}
Compare a 30-day training run on 512 H100s:
\textbf{Poland} vs \textbf{Quebec}. How many tonnes of CO$_2$ saved?
\end{alertblock}

\begin{lstlisting}
fleet = mlsysim.Fleet(
    name="Training Fleet",
    node=mlsysim.Systems.Nodes.DGX_H100,
    count=64,
    fabric=mlsysim.Systems.Fabrics.InfiniBand_NDR)
solver = mlsysim.SustainabilityModel()
# YOUR CODE: solve for Poland and Quebec
\end{lstlisting}
\end{frame}

% --- 5.9b Exercise Part B: Multi-Vendor TCO Shootout ---
\begin{frame}[fragile]{Exercise 5b: Multi-Vendor TCO Shootout}
\note{[5 min] THE ISCA QUESTION. Which vendor gives the cheapest training?
Expected: MI300X cluster is cheapest per FLOP (192 GB HBM means fewer
nodes needed for large models), but H100 has the best software ecosystem.
Gaudi\,3 is competitive on raw cost. The answer depends on model size.
Turn to your neighbor: did you get the same answer? Why or why not? 60 seconds.}

\begin{alertblock}{Your Task (5 minutes)}
Same workload: Llama-3 70B, 30-day training run, 512 accelerators.\\
Which cluster is cheapest for the \textbf{same throughput}?
\end{alertblock}

\begin{lstlisting}
econ = mlsysim.EconomicsModel()
for cluster_name in ["H100", "MI300X", "Gaudi3"]:
    hw = getattr(mlsysim.Hardware.Cloud, cluster_name)
    fleet = mlsysim.Fleet(name=cluster_name, hw=hw,
                          count=64)
    r = econ.solve(fleet=fleet, duration_days=30,
                   infrastructure_multiplier=2.5)
    print(f"{cluster_name:>8}: TCO=${r.tco_usd:,.0f}")
\end{lstlisting}

\small
\alert{The cheapest accelerator depends on the workload's binding constraint.}
\end{frame}

% --- 5.10 Takeaway ---
\begin{frame}{Key Takeaway: Economics \& Sustainability}
\note{[1 min] Three points, repeat each.}
\begin{center}
\Large
\begin{enumerate}
  \item \textbf{TCO $\neq$ GPU cost.} Infrastructure is 2--2.5$\times$.
  \item \textbf{Utilization is the cost lever.} Idle GPUs cost the same.
  \item \textbf{Geography $>$ algorithms} for carbon reduction.
\end{enumerate}

\vspace{1em}
\normalsize
\textit{``The cheapest watt is the one you don't consume.\\
The cleanest watt is the one from a hydro dam.''}
\end{center}
\end{frame}


% =====================================================================
% PART 6: DESIGN SPACE EXPLORATION
% =====================================================================


\begin{frame}[fragile]{Cross-Domain Carbon Accounting (Scenario B)}
Training time depends on hardware efficiency. Carbon footprint depends on geographic grid intensity. \mlsysim{} composes solvers cleanly.
\begin{lstlisting}[language=Python]
from mlsysim.solvers import DistributedModel, SustainabilityModel
from mlsysim.systems.registry import Systems
from mlsysim.infra.registry import Infrastructure
from mlsysim.models.registry import Models

fleet = Systems.Clusters.Research_256  # 32x DGX H100 nodes

perf = DistributedModel().solve(
    model=Models.Language.Llama3_70B, fleet=fleet,
    tp_size=8, pp_size=1, dp_size=32)
training_days = perf.latency.to("days").magnitude

sust = SustainabilityModel().solve(
    fleet=fleet, duration_days=training_days,
    datacenter=Infrastructure.Datacenters.GCP_Iowa,  # High-carbon grid
    mfu=perf.mfu)
print(f"Total Carbon: {sust.carbon_footprint:.1f} tonnes CO2e")
\end{lstlisting}
\end{frame}


\end{document}
